% Design, Fig. sectors (a): a parity that a gate conserves. Bits of a string: x and z per qubit (3 qubits here).
% The parity of a term under the row r is the sum of its bits in the columns where r has a 1 (teal bands), modulo 2.
% The rotation about Q branches P into P and the new term P xor Q; since r.Q = 0, both have the parity of P and
% stay in the same sector.
\begin{tikzpicture}[x=1cm, y=1cm]
  \def\cw{0.40}\def\x0{2.80}\def\rh{0.46}
  \newcommand{\bits}[4]{% row, six bits, fill, extra style
    \foreach \b [count=\i from 0] in {#2}
      \node[cl, fill=#3, minimum width=\cw cm, minimum height=0.32cm, anchor=center, #4]
        at ({\x0+\cw*\i+(\i>2)*0.12}, {-0.62-\rh*#1}) {\b};}
  % the columns where the row has a 1
  \foreach \i in {0,1,4}
    \fill[rowc!45] ({\x0+\cw*\i+(\i>2)*0.12-\cw/2}, -0.28) rectangle ++(\cw, {-0.30-\rh*3-0.26});
  \foreach \t [count=\i from 0] in {$x_1$,$x_2$,$x_3$,$z_1$,$z_2$,$z_3$}
    \node[stxt] at ({\x0+\cw*\i+(\i>2)*0.12}, -0.16) {\t};
  \node[stxt, anchor=east] at (\x0+3.05, -0.16) {};
  % rows
  \node[txt, anchor=west] at (0,-0.62) {Row $r$};
  \bits{0}{1,1,0,0,1,0}{rowc}{}
  \node[txt, anchor=west] at (0,{-0.62-\rh}) {Gate $Q$ (IYI)};
  \bits{1}{0,1,0,0,1,0}{genc}{}
  \node[txt, anchor=west] at (0,{-0.62-2*\rh}) {Term $P$ (XZX)};
  \bits{2}{1,0,1,0,1,0}{white}{}
  \node[txt, anchor=west] at (0,{-0.62-3*\rh}) {New $P\oplus Q$ (XXX)};
  \bits{3}{1,1,1,0,0,0}{white}{draw=newc, line width=1.0pt}
  % parities
  \node[stxt] at (5.90,-0.16) {Parity};
  \foreach \v [count=\r from 1] in {0, 0, 0}
    \node[txt] at (5.90, {-0.62-\rh*\r}) {\v};
  \draw[decorate, decoration={brace, amplitude=2.5pt}, line width=0.6pt] (6.17,{-0.62-2*\rh+0.16}) -- (6.17,{-0.62-3*\rh-0.16});
  \node[txt, anchor=west, align=left] at (6.27,{-0.62-2.5*\rh}) {Same\\Sector};
  % the rule
  \node[key] at (3.75,-2.55) {$r\cdot Q = 0 \;\Rightarrow\; r\cdot(P\oplus Q) = r\cdot P$};
\end{tikzpicture}
